% Part: incompleteness
% Chapter: incompleteness-provability
% Section: godels-paper

\documentclass[../../../include/open-logic-section]{subfiles}

\begin{document}

\olfileid{inc}{inp}{gop}

\olsection{గ్యోడెల్ అసలు పరిశోధనపత్రంతో పోలిక}

గ్యోడెల్ 1931 పరిశోధనపత్రాన్ని కొంత సమయం వెచ్చించి చదవడం విలువైనది. దాని
పరిచయం మనం ఇప్పుడే చర్చించిన ఆలోచనలను స్థూలంగా వివరిస్తుంది. ఆ పత్రాన్ని
పైపైన చదివినా ప్రతి దశలో ఏమి జరుగుతోందో సులభంగా తెలుస్తుంది: మొదట గ్యోడెల్
అధికారిక వ్యవస్థ $P$ను (వాక్యనిర్మాణం, స్వీకృతాలు, నిరూపణ నియమాలు)
వివరిస్తాడు; తర్వాత ఆదిమ పునరావర్తన ప్రమేయాలను, సంబంధాలను నిర్వచిస్తాడు;
తర్వాత $x B y$ ఆదిమ పునరావర్తన సంబంధమని చూపి, ఆదిమ పునరావర్తన ప్రమేయాలు,
సంబంధాలు $\Th{P}$లో ప్రాతినిధ్యం చేయబడతాయని వాదిస్తాడు. ఆపై పైన చెప్పిన
విధంగానే అసంపూర్ణతా సిద్ధాంతాన్ని నిరూపిస్తాడు. విభాగం~3లో,
వ్యుత్పాదించలేని ప్రకటనను అంకగణిత భాషలోని వాక్యంగా తీసుకోవచ్చని చూపిస్తాడు.
ఇదే $\beta$-ఉపసిద్ధాంతానికి మూలం; పునరావర్తన ప్రమేయాలు $\Th{Q}$లో
ప్రాతినిధ్యం చేయబడతాయని చూపేటప్పుడు వరుసలను నిర్వహించేందుకు మనమూ దీనినే
వాడాం. సరిపోయే కనిష్ఠ స్వీకృతాల సమితిని వేరుగా గుర్తించేంతవరకు గ్యోడెల్
వెళ్లలేదు; కానీ ఇప్పుడు $\Th{Q}$ ఆ పనికి సరిపోతుందని మనకు తెలుసు. చివరగా,
విభాగం~4లో రెండవ అసంపూర్ణతా సిద్ధాంతపు నిరూపణను స్థూలంగా వివరిస్తాడు.

\end{document}
